\section{Recovery sets and normalized equations}
\label{sec:recovery-equations}

Let $C\leq\F_q^E$ be a linear code.  We first separate four layers that are
often harmlessly conflated when only locality is at issue.

Fix a target $x\in E$ and a helper bound $r$.  A dual word $w\in C^\perp$
with $w_x\ne0$ gives
\[
 c_x=\sum_{y\ne x}-\frac{w_y}{w_x}c_y
 \qquad(c\in C).
\]
Its exact helper support is $\supp(w)\setminus\{x\}$.  Define
\[
 \cH_x^{(\le r)}(C)=
 \{\supp(w)\setminus\{x\}:w\in C^\perp, w_x\ne0,
       \wt(w)\le r+1\}.
\]
The usual bounded recovery-set family is the upward closure of
$\cH_x^{(\le r)}(C)$ inside the subsets of $E\setminus\{x\}$ of size at most
$r$.  Its inclusion-minimal members form the minimal recovery-set clutter.
Thus exact supports, recovery sets, and minimal recovery sets are distinct
but determine the same monotone success event.

The coefficient layer is
\[
 \cP_x^{(\le r)}(C)=
 \{w\in C^\perp:w_x=1, \wt(w)\le r+1\}.
\]
These normalized recovery equations determine their exact supports, while
the converse generally fails.  The operational model is scalar exact repair:
each participating helper supplies one symbol of $\F_q$.  Array-code access,
subpacketization, and repair bandwidth are different optimization problems
\cite{ShahEtAl2010,YeBarg2016,KurzYaakobi2021}.

If $A\subseteq E\setminus\{x\}$ is the surviving helper set, put
\[
 f_{x,r}(A)=1
 \quad\Longleftrightarrow\quad
 \text{$A$ contains a member of $\cH_x^{(\le r)}(C)$}.
\]
Under independent helper survival with probabilities $(s_y)$, the bounded
repair reliability is
\[
 R_{x,r}((s_y))=\mathbb E[f_{x,r}].
\]
It depends only on the minimal recovery sets, but it is not determined by
their minimum cardinality.  Section~\ref{sec:separations} shows that even the
complete relative-weight hierarchy does not determine the corresponding
bounded reliability law.

For several target coordinates, a recovery equation is a dual subspace rather
than one dual word.  Let $P\subseteq E$, $J=E\setminus P$, and let
$G=(G_P\mid G_J)$ be a full-row-rank generator matrix.  If
$T\leq\operatorname{im}G_P$ is a subspace of target linear combinations, a
normalized recovery system for $T$ consists of linear maps
\[
 \alpha:T\longrightarrow\F_q^P,
 \qquad
 \beta:T\longrightarrow\F_q^J
\]
such that
\[
 G_P\alpha=\operatorname{id}_T,
 \qquad
 G_J\beta=-\operatorname{id}_T.
\]
For each $t\in T$, the vector $(\alpha(t),\beta(t))$ is then a dual word.
The map $\alpha$ records the chosen target normalization and $\beta$ records
the helper coefficients.

The helper union of the system is the support of $\beta(T)$:
\[
 \supp(\beta(T))=
 \{j\in J:\text{some }t\in T\text{ has }\beta(t)_j\ne0\}.
\]
Its helper cost is $|\supp(\beta(T))|$.  We call a target subspace internally
recoverable when it lies in
\[
 W_P=\operatorname{im}G_P\cap\operatorname{im}G_J.
\]
Only this recovered subspace dimension is counted below.  Relations supported
entirely on $P$ have recovered dimension zero and do not consume helpers.

For later comparison, the normalized equations at one coordinate reconstruct
the represented code at radius $r$ when their span is $C^\perp$.  The least
such $r$ is the reconstruction radius.

\begin{theorem}[MDS reconstruction from minimum-weight recovery equations]
\label{thm:mds-reconstruction}
\coverage{absent}%
Let $C\leq\F_q^E$ be an $[n,k]_q$ MDS code with $1\leq k<n$, and let
$x\in E$.  Then
\[
 \operatorname{span}\cP_x^{(\leq k)}(C)=C^\perp,
\]
the reconstruction radius is $k$, and the supports of the minimum-weight
recovery equations form the complete $k$-uniform clutter on
$E\setminus\{x\}$.
\end{theorem}

\begin{proof}
The dual has dimension $n-k$ and minimum distance $k+1$.  For every
$(k+1)$-set $T\subseteq E$, restriction of $C^\perp$ to $E\setminus T$ has
domain dimension $n-k$ and codomain dimension $n-k-1$.  Its kernel contains a
nonzero word supported on $T$, and the distance bound forces its support to be
all of $T$.  Hence the exact helper supports through $x$ are precisely the
$k$-subsets of $E\setminus\{x\}$.

Choose a fixed $(k-1)$-set $Q\subseteq E\setminus\{x\}$.  For every
$t\in E\setminus(\{x\}\cup Q)$, choose the normalized minimum dual word
supported on $\{x,t\}\cup Q$.  Evaluation at the private coordinates $t$
shows that these $n-k$ words are independent.  They are therefore a basis of
$C^\perp$.  No smaller radius contains a nonzero normalized equation, again
by the dual-distance bound.
\proves{thm:mds-reconstruction}%
\end{proof}

The support clutter in this theorem depends only on the MDS parameters.  The
normalized equations retain the represented dual code.  This is the first
instance of the support--coefficient distinction used in the transfer and
separation results below.
